\documentclass[10pt,a4paper]{amsart}
\usepackage[margin=19mm]{geometry}
\usepackage{amsmath,amssymb,mathtools}
\usepackage[scr=rsfs,cal=euler]{mathalfa}
\usepackage{xfrac}
\usepackage[unicode,hidelinks,bookmarks=false]{hyperref}
\newcommand{\ie}{i.e.\ }
\newcommand{\cf}{cf.\ }
\newcommand{\resp}{respectively\ }
\newcommand{\Subsection}{\subsection}
\providecommand{\nobreakdash}{\nobreak\hbox{-}}
\setlength{\emergencystretch}{2em}
\multlinegap0pt
\usepackage{amssymb}
\usepackage[shortlabels]{enumitem}
\usepackage{tikz}
\newtheorem{theorem}{Theorem}
\newtheorem{lemma}{Lemma}
\newcommand{\A}{\mathcal A}
\newcommand{\conf}{\operatorname{conf}}
\title{Combinatorics of Inflection Points of Plane Curve Shadows}
\author{Boris Shapiro}
\date{Source: arXiv:2605.27471v1 (CC BY 4.0)}
\begin{document}
\maketitle

\noindent\emph{Source and licence.} Combinatorics of Inflection Points of Plane Curve Shadows. Version arXiv:2605.27471v1 (CC BY 4.0), distributed by arXiv under the Creative Commons Attribution 4.0 International licence, http://creativecommons.org/licenses/by/4.0/, as recorded in the arXiv licence metadata for this exact version and shown on the abstract page https://arxiv.org/abs/2605.27471v1. Full licence notice: \texttt{licenses/arxiv-2605.27471-CC-BY-4.0.txt}.\par\smallskip

\noindent\textit{Editorial scope.} Counted unit: the article's theorem thm:formula (source0.tex lines 254--261). The article's complete original proof of it is delivered with it (source0.tex lines 263--272, one \texttt{proof} environment of 10 lines). No mathematics is written, repaired or re-derived: the statement and the proof are the article's own lines, in the article's own notation, and no retained line is substituted or re-flowed.  Retained uncounted as the source-local context the proof needs: lines 204--206; lines 221--223.  Nothing was omitted from any retained span, so the package carries no omission record.  No retained line contains a citation, so the wrapper carries no bibliography and no bibliography entry.  No retained line contains a figure environment, an image-inclusion command or drawing code, so no image is delivered and the package declares no asset dependency.  Editorial formatting only, in addition: an AMS \texttt{amsart} preamble that restates the article's own declarations and macros used by the retained lines, namely the environments \texttt{theorem}, \texttt{lemma} on separate counters, together with the standard packages those lines need.  The retained environments print in this document as Theorem 1, Lemma 1 in order of appearance, whereas the article prints the same items with the numbering of its own full document.  The complete native source of the article as distributed in the arXiv e-print for this exact version is bundled unchanged as \texttt{sources/201490/source0.tex}.
\begin{theorem}[inflectionless criterion]\label{thm:shapiro}
Let $c$ be a tree-like generic plane curve.  Then $c$ is carried by an orientation-preserving diffeomorphism of the plane to a curve without inflection points if and only if its decomposition admits an admissible local coorientation.
\end{theorem}

\begin{lemma}[Normalization]\label{lem:normalization}
Every generic curve is isotopic, through generic curves, to a normalized one without increasing the number of inflections.
\end{lemma}

\begin{theorem}\label{thm:formula}
Let $c$ be a tree-like generic plane curve.  Then
\[
 \mu([c])=
 \min_{C\in\A(c)} \conf(C),
\]
where $\A(c)$ is the finite set of admissible local coorientations of the building polygons of $c$.
\end{theorem}

\begin{proof}
First let $c'$ be a normalized representative of $[c]$.  On every side away from the vertex neighborhoods, the curve has a definite convex side.  Hence $c'$ determines a local coorientation $C$.  The necessity part of Theorem~\ref{thm:shapiro}, applied locally to the inflection-free sides of the decomposition, shows that $C$ is admissible.  Each normalized inflection occurs exactly at a vertex at which the two induced convex sides are incompatible.  Therefore
\[
 \#\{\text{inflections of }c'\}\ge \conf(C)\ge
 \min_{D\in\A(c)}\conf(D).
\]
Taking the minimum over normalized representatives and using Lemma~\ref{lem:normalization} gives the lower bound.

Conversely, take an admissible local coorientation $C$.  The constructive proof of Theorem~\ref{thm:shapiro} realizes all building polygons with the prescribed convex sides.  At every non-conflicting gluing one smooths the corner with curvature of the prescribed sign.  At every conflict, one inserts a single normalized inflection in a sufficiently small neighborhood of the corresponding double point.  Since the block-adjacency graph is a tree, these local choices are independent.  The resulting curve is isotopic to $c$ and has exactly $\conf(C)$ inflection points.  Minimizing over $C$ proves the reverse inequality.
\end{proof}

\end{document}
